\section{数据预处理}

\subsection{数据清洗}

在问卷收集完成并汇总后，首先对问卷进行系统性审核，剔除无效问卷，仅保留有效样本开展后续分析，以保障数据质量与研究有效性。无效问卷的剔除标准如下：

\textbf{（1）}纸质问卷中，字迹模糊、问卷有效内容破损，导致无法提取有效信息的问卷，视为无效问卷并剔除。

\textbf{（2）}电子问卷中，答题总时长不超过50秒的问卷，视为无效问卷并剔除。

\textbf{（3）}所有问卷中，出现连续3题及以上未作答、答案与问题无关、答案数量与题目要求不一致、答卷中大部分内容选择同一选项，或出现明显规律性作答痕迹的问卷，均视为无效问卷并剔除。

\subsection{缺失值处理}

对于清洗后有效问卷中仍存在的缺失值，本研究采用如下处理方式：

首先对缺失率超过10\%的变量进行评估，主要是“居家办公情况”一题，约有36.62\%的受访者为学生、退休人员或无业待业者，缺失比例较高（如图\ref{fig:homeoffice}），将不进行插补，并在描述性统计中单独进行分析。除此之外，对于“不清楚居所是否为隔音墙”选项，将其与“无隔音墙”合并为一类处理，后续分析中不再单独区分。

对于缺失率低于10\%的变量，分类变量采用该变量的众数进行填充，连续变量（如量表题）采用均值或中位数进行填充，以在不显著影响样本量的前提下，降低缺失值对后续统计分析的干扰。

\fig{图片9.png}{居家办公题项缺失情况\label{fig:homeoffice}}

\subsection{数据编码}

为实现问卷数据的标准化处理与后续量化分析，本研究对回收后的有效问卷进行统一的数据编码处理，将问卷中各类文字、选项信息转化为可用于统计分析的数字代码。

针对性别、年龄、居住地、职业、文化程度等人口统计学类单选题，按照选项顺序依次赋予对应数字编码实现量化转换。对于受噪声影响程度、隔音水平评价、治理满意度等李克特五级量表类题目，将程度表述依次赋值为1-5分，统一量化标准以便开展均值、方差及相关性分析。针对多选类题目，采用二分法编码形式，将受访者选择的选项标记为1，未选择项标记为0，精准区分不同噪声类型、应对方式等的选择差异。针对开放题与访谈文本内容，采用质性编码方式，围绕噪声干扰类型、健康影响、治理诉求、政策认知等核心主题进行概念提炼，归纳形成核心范畴，将文字性访谈资料转化为可分析的定性数据。具体如下表：

\begin{table}[H]
\centering
\caption{问卷编码规则}
\resetrowcolors
\begin{tabular}{M{2.4cm}M{3.4cm}M{2.4cm}p{5.6cm}}
\bluetoprule
\rowcolor{tabhead}
\thead{变量分类} & \thead{变量名称} & \thead{变量类型} & \multicolumn{1}{c}{\thead{编码规则}} \\
\hline
\multirow{8}{*}{基本信息} & 性别 & 二分类变量 & 1=男，2=女 \\ \cline{2-4}
 & 年龄 & 有序分类变量 & 1=第一个选项，2=第二个选项，以此类推 \\ \cline{2-4}
 & 居住地 & 无序分类变量 & 1=第一个选项，2=第二个选项，以此类推 \\ \cline{2-4}
 & 居住区域类型 & 无序分类变量 & 1=第一个选项，2=第二个选项，以此类推 \\ \cline{2-4}
 & 职业类型 & 无序分类变量 & 1=第一个选项，2=第二个选项，以此类推 \\ \cline{2-4}
 & 文化程度 & 有序分类变量 & 1=第一个选项，2=第二个选项，以此类推 \\ \cline{2-4}
 & 是否为隔音墙 & 无序分类变量 & 1=第一个选项，2=第二个选项，以此类推 \\ \cline{2-4}
 & 居家办公情况 & 有序分类变量 & 1=第一个选项，2=第二个选项，以此类推 \\ \hline
\multirow{6}{*}{噪声暴露与环境感知} & 受噪声影响程度 & 李克特5级量表 & 1=影响非常大，5=几乎没有影响 \\ \cline{2-4}
 & 隔音水平评价 & 李克特5级量表 & 1=很好，5=很差 \\ \cline{2-4}
 & 各类噪声干扰程度 & 李克特5级量表 & 1=无干扰，5=严重干扰 \\ \cline{2-4}
 & 离噪声源距离 & 有序分类变量 & 1=第一个选项，2=第二个选项，以此类推 \\ \cline{2-4}
 & 噪音影响时段 & 有序分类变量 & 1=第一个选项，2=第二个选项，以此类推 \\ \cline{2-4}
 & 噪声负面影响 & 李克特5级量表 & 1=不认同，5=认同 \\ \hline
\multirow{3}{*}{噪声治理满意度} & 声环境质量满意度 & 李克特5级量表 & 1=非常不满意，5=非常满意 \\ \cline{2-4}
 & 噪声治理满意度 & 李克特5级量表 & 1=非常不满意，5=非常满意 \\ \cline{2-4}
 & 主要应对方式 & 多重响应变量 & 采用二进制编码，被选择编码为1，未选择编码为0 \\ \hline
\multirow{4}{*}{噪声治理认知} & 投诉渠道认知 & 李克特5级量表 & 1=非常不了解，5=非常了解 \\ \cline{2-4}
 & 国家政策认知 & 李克特5级量表 & 1=非常不了解，5=非常了解 \\ \cline{2-4}
 & 治理必要性认知 & 李克特5级量表 & 1=完全没有必要，5=非常有必要 \\ \cline{2-4}
 & 责任主体认知 & 多重响应变量 & 采用二进制编码，被选择编码为1，未选择编码为0 \\ \hline
\multirow{2}{*}{噪声治理建议} & 举措优先度 & 多重响应变量 & 采用二进制编码，被选择编码为1，未选择编码为0 \\ \cline{2-4}
 & 具体意见 & 无结构文本数据 & 开放式编码 \\
\bluebottomrule
\end{tabular}
\end{table}

